\documentclass[../../main.tex]{subfiles}
\begin{document}

\begin{flushright}
\footnotesize\textit{【自動文字起こし・要確認】}
\end{flushright}


{\bf [解]}
\paragraph{(1)}

与えられた数列
\begin{align}
 I(n) = \int_{0}^{\frac{n\pi}{2}} |\sin x|\, dx \label{eq:1}
\end{align}
について考える．自然数$k$に対して新しく
\begin{align}
 a_k = \int_{(k-1)\pi/2}^{k\pi/2}|\sin x|\,dx
\end{align}
とおく．$|\sin x|$は周期$\pi$を持ち，この区間での値は$[0,\pi/2]$または$[\pi/2,\pi]$での$\sin t$の値に一致するので
\begin{align}
 a_k=\int_0^{\pi/2} \sin t\,dt=1 \label{eq:ak}
\end{align}
が成り立つ．

従って\cref{eq:1}は
\begin{align}
 I(n) = \sum_{k=1}^{n} a_k = n \label{eq:2}
\end{align}
である．

\paragraph{(2)}

$\displaystyle\int_0^{\frac{s\pi}{2}}\cos x\,dx=\sin\dfrac{s\pi}{2}$であるから，示すべき不等式は
\begin{align}
 s\le\sin\dfrac{s\pi}{2}\le\dfrac{\pi}{2}s\qquad(0\le s\le1) \label{eq:3}
\end{align}
と同値である．左右の不等号を別々に示す．

\subparagraph{左側の不等号}

$f(s)=\sin\dfrac{s\pi}{2}-s$とおいて$f(s)\ge0$を示す．
区間の両端での値は$f(0)=0$，$f(1)=\sin\dfrac{\pi}{2}-1=0$であり，一階微分は
\begin{align}
f'(s)=\dfrac{\pi}{2}\cos\dfrac{s\pi}{2}-1
\end{align}
となる．これはは単調減少であり，$\cos\dfrac{\alpha\pi}{2}=\dfrac{2}{\pi}$をみたす$\alpha\in(0,1)$を境に符号が$+$から$-$に変わる．
よって$f(s)$の増減表は以下のようになる．

従って区間$[0,1]$全体で$f(s)\ge0$，すなわち$s\le\sin\dfrac{s\pi}{2}$であり\cref{eq:3}の左側の不等式は示された．

\subparagraph{右側の不等号}

$h(s)=\dfrac{\pi}{2}s-\sin\dfrac{s\pi}{2}$とおいて$h(s) \ge 0$を示す．
一階微分は
\begin{align}
h'(s)=\dfrac{\pi}{2}\Bigl(1-\cos\dfrac{s\pi}{2}\Bigr)\ge0
\end{align}
と非負だから$h(s)$は単調増加であり，よって区間$[0,1]$では
\begin{align}
 h(s)\ge h(0)=0
\end{align}
である．従って$\sin\dfrac{s\pi}{2}\le\dfrac{\pi}{2}s$であり右側の不等号は示された．
\casesend

以上で\cref{eq:3}，従って題意は示された．


\paragraph{(3)}

まず題意の不等式の積分において$at\to t$と置換すると
\begin{align}
 0 \le \dfrac{1}{a}\int_{0}^{a\pi/2}|\sin t| dt - 1 \le \left(\dfrac{\pi}{2}-1\right)\left(1-\dfrac{[a]}{a}\right) \\
 a \le \int_{0}^{a\pi/2}|\sin t| dt \le \left(\dfrac{\pi}{2}-1\right)\left(a-[a]\right)+a \label{eq:4}
\end{align}
であるからこれを示せば良い．

題意より$[a]$が奇数なので
\begin{align}
 [a]=2m-1 \quad (m\in\mathbb N) \label{eq:5}
\end{align}
と表せる．ガウス記号の性質より
\begin{align}
 \beta=a-[a]\in[0,1) \label{eq:6}
\end{align}
とおける．

 \cref{eq:4}の積分を
\begin{align}
 J(a)= \int_0^{\frac{a\pi}{2}}|\sin t|\,dt
\end{align}
とおいて，これを変形すると\cref{eq:2}を用いて
\begin{align}
 J(a)
  &=\int_0^{\frac{(2m-1)\pi}{2}}|\sin t|\,dt+\int_{\frac{(2m-1)\pi}{2}}^{\frac{a\pi}{2}}|\sin t|\,dt\\
  &=I(2m-1) + \int_{\frac{(2m-1)\pi}{2}}^{\frac{a\pi}{2}}|\sin t|\,dt\\
  &=(2m-1) + \int_{\frac{(2m-1)\pi}{2}}^{\frac{a\pi}{2}}|\sin t|\,dt
\end{align}
となる．(1)と同様$\sin x$の周期性より$m\pi$だけ積分区間を移動してよく，
\begin{align}
 J(a)
  &=(2m-1) + \int_{\frac{-\pi}{2}}^{\frac{\beta \pi}{2}-\frac{\pi}{2}}|\sin t|\,dt \\
  &=(2m-1) + \int_{0}^{\frac{\beta \pi}{2}}\left|\sin \left( t - \dfrac{\pi}{2}\right)\right|\,dt  \\
  &=(2m-1) + \int_{0}^{\frac{\beta \pi}{2}}\left|\cos t \right|\,dt \\
  &=(2m-1) + \int_{0}^{\frac{\beta \pi}{2}}\cos t \,dt 
\end{align}
である．\cref{eq:6}より$0 \le \beta < 1$だからここに(2)の不等式を代入して上下から評価すると
\begin{align}
 (2m-1) + \beta \le J(a) \le (2m-1) + \left(\dfrac{\pi}{2}-1\right) \beta \\{}
 [a] + \beta \le J(a) \le [a] + \left(\dfrac{\pi}{2}-1\right) \beta \\
\end{align}
である．\cref{eq:6}を代入して
\begin{align}
 a \le J(a) \le [a] + \left(\dfrac{\pi}{2}-1\right)\left(a-[a]\right) \label{eq:7}
\end{align}
を得る．これは\cref{eq:4}に他ならないから，\cref{eq:4}が示された．
よって題意は示された．


{\bf [解説]}

(2)の途中で出てきた不等式はジョルダンの不等式と呼ばれる．図形的に示すとわかりやすい．
% TODO :: ジョルダンの不等式の図形

\newpage

\end{document}
